\chapter{1958年湖南卷}

\section{解答题}

\begin{problem}
不用除法证明 \(\left(2+x\right)^{2n-1}+x^{2n-1}+x+1\) 能被 \(x+1\) 整除（\(n\) 为正整数）.
\end{problem}

\begin{solution}
令 \(f(x)=\left(2+x\right)^{2n-1}+x^{2n-1}+x+1\). 因为 \(n\) 为正整数，所以 \(2n-1\) 是奇数. 代入 \(x=-1\)，得
\[
f(-1)=\left(2-1\right)^{2n-1}+\left(-1\right)^{2n-1}-1+1=1-1-1+1=0
\]
根据因式定理，多项式 \(f(x)\) 在 \(x=-1\) 时的值为零，等价于一次式 \(x-(-1)=x+1\) 是 \(f(x)\) 的因式. 因此原式能被 \(x+1\) 整除，且上述证明没有使用多项式除法.
\end{solution}

\begin{problem}
不解方程 \(6x^{2}-x-2=0\)，求作一个二次方程，使它的根是原方程的根的倒数.
\end{problem}

\begin{solution}
设原方程的两个根为 \(x_{1}\)，\(x_{2}\). 由根与系数的关系，
\(x_{1}+x_{2}=\frac{1}{6}\)，\(x_{1}x_{2}=-\frac{1}{3}\).
因为 \(x_{1}x_{2}\ne0\)，所以两个根均不为零. 令它们的倒数分别为 \(y_{1}=\frac{1}{x_{1}}\)，\(y_{2}=\frac{1}{x_{2}}\)，则
\(y_{1}+y_{2}=\frac{x_{1}+x_{2}}{x_{1}x_{2}}=\frac{\frac{1}{6}}{-\frac{1}{3}}=-\frac{1}{2}\)，\(y_{1}y_{2}=\frac{1}{x_{1}x_{2}}=\frac{1}{-\frac{1}{3}}=-3\).
以 \(y_{1}\)，\(y_{2}\) 为根的二次方程应为
\[
y^{2}-\left(y_{1}+y_{2}\right)y+y_{1}y_{2}=0
\]
代入上面求出的和与积，得
\[
y^{2}+\frac{1}{2}y-3=0
\]
将方程两边同乘以 \(2\)，再把未知数改记为 \(x\)，得到所求方程
\[
2x^{2}+x-6=0
\]
\end{solution}

\begin{problem}
求证 \(\sin75^{\circ}=\frac{1}{4}\left(\sqrt{6}+\sqrt{2}\right)\).
\end{problem}

\begin{solution}
由正弦的和角公式，
\[
\sin75^{\circ}=\sin\left(45^{\circ}+30^{\circ}\right)=\sin45^{\circ}\cos30^{\circ}+\cos45^{\circ}\sin30^{\circ}
\]
又因为 \(\sin45^{\circ}=\cos45^{\circ}=\frac{\sqrt{2}}{2}\)，\(\cos30^{\circ}=\frac{\sqrt{3}}{2}\)，\(\sin30^{\circ}=\frac{1}{2}\)，所以
\[
\sin75^{\circ}=\frac{\sqrt{2}}{2}\cdot\frac{\sqrt{3}}{2}+\frac{\sqrt{2}}{2}\cdot\frac{1}{2}=\frac{\sqrt{6}}{4}+\frac{\sqrt{2}}{4}=\frac{1}{4}\left(\sqrt{6}+\sqrt{2}\right)
\]
故等式成立.
\end{solution}

\begin{problem}
证明：设三角形的三个角成等差数列，最大角与最小角相差 \(30^{\circ}\)，则其三边之比顺次等于 \(2\sqrt{2}\)，\(2\sqrt{3}\)，\(\sqrt{6}+\sqrt{2}\) 之比.
\end{problem}

\begin{solution}
将三角形的三个角按从小到大记为 \(A\)，\(B\)，\(C\). 因为三个角成等差数列，所以 \(A+C=2B\). 又因为 \(A+B+C=180^{\circ}\)，所以 \(3B=180^{\circ}\)，即 \(B=60^{\circ}\). 由最大角与最小角相差 \(30^{\circ}\)，得 \(C-A=30^{\circ}\). 因而
\(A+C=180^{\circ}-B=120^{\circ}\)，\(C-A=30^{\circ}\).
两式相加、相减，得到 \(C=75^{\circ}\)，\(A=45^{\circ}\).

设 \(a\)，\(b\)，\(c\) 分别为角 \(A\)，\(B\)，\(C\) 所对的边. 由正弦定理，
\[
a:b:c=\sin45^{\circ}:\sin60^{\circ}:\sin75^{\circ}
\]
其中 \(\sin45^{\circ}=\frac{\sqrt{2}}{2}\)，\(\sin60^{\circ}=\frac{\sqrt{3}}{2}\)，并且由和角公式
\[
\sin75^{\circ}=\sin\left(45^{\circ}+30^{\circ}\right)=\frac{\sqrt{2}}{2}\cdot\frac{\sqrt{3}}{2}+\frac{\sqrt{2}}{2}\cdot\frac{1}{2}=\frac{\sqrt{6}+\sqrt{2}}{4}
\]
所以
\[
a:b:c=\frac{\sqrt{2}}{2}:\frac{\sqrt{3}}{2}:\frac{\sqrt{6}+\sqrt{2}}{4}=2\sqrt{2}:2\sqrt{3}:\left(\sqrt{6}+\sqrt{2}\right)
\]
这正是题目所要求的三边之比.

\solutionfigure{\bitmapfigure[max width=0.42\linewidth]{img/1958/hunan/q4_solution_fig1.png}}
\end{solution}

\begin{problem}
一矩形内接于圆，它的面积是 \(48\mathrm{~cm}^{2}\)，圆半径是 \(5\mathrm{~cm}\)，求矩形的两邻边.
\end{problem}

\begin{solution}
\solutionfigure{\bitmapfigure[max width=0.42\linewidth]{img/1958/hunan/q5_solution_fig1.png}}

设矩形的两邻边为 \(AB=x\)，\(BC=y\)，其中 \(x>0\)，\(y>0\). 矩形的两条对角线相等且互相平分，所以外接圆的圆心是两条对角线的交点，\(AC\) 是圆的直径. 因此
\[
AC=2\times5\mathrm{~cm}=10\mathrm{~cm}
\]
由矩形面积和勾股定理，分别得到
\(xy=48\)，\(x^{2}+y^{2}=AC^{2}=100\).
于是
\[
(x+y)^{2}=x^{2}+2xy+y^{2}=100+2\times48=196
\]
因为 \(x>0\)，\(y>0\)，所以 \(x+y>0\)，从而 \(x+y=14\). 因此 \(x\)，\(y\) 是方程
\[
t^{2}-14t+48=0
\]
的两个根. 因式分解得
\[
(t-6)(t-8)=0
\]
故 \(\{x,y\}=\{6,8\}\). 所以矩形的两邻边为 \(6\mathrm{~cm}\) 和 \(8\mathrm{~cm}\)，顺序可以互换.
\end{solution}

\begin{problem}
设 \(AD\)，\(BE\) 是由圆 \(O\) 的直径 \(AB\) 两端点所引的切线，\(DE\) 是过圆上任意一点 \(F\) 的切线，此切线与 \(AD\)，\(BE\) 交于 \(D\) 及 \(E\)，求证：\(AB\cdot BE=OF^{2}\).
\end{problem}

\begin{solution}
\solutionfigure{\bitmapfigure[max width=0.42\linewidth]{img/1958/hunan/q6_solution_fig1.png}}

先核对题面中的等式. 原件的题干印作 \(AB\cdot BE=OF^{2}\)，但同页原解答的末行印作 \(AD\cdot BE=OF^{2}\). 前一个等式按题面构型并不恒成立：取 \(F\) 为与直径 \(AB\) 垂直的直径端点，记 \(OF=r\). 建立直角坐标系，使 \(O\) 为原点，\(OF\) 为负的横轴方向，则可取
\(A=(0,r)\)，\(B=(0,-r)\)，\(F=(-r,0)\).
三点处的切线分别为 \(y=r\)，\(y=-r\)，\(x=-r\)，所以 \(D=(-r,r)\)，\(E=(-r,-r)\). 从而 \(AD=BE=OF=r\)，而 \(AB=2r\)，所以
\[
AB\cdot BE=2r^{2}\ne r^{2}=OF^{2}
\]
因此原题的题干与原解答存在客观冲突，不能把题干中的 \(AB\cdot BE=OF^{2}\) 当作恒成立命题. 下面完整证明原解答末行给出的正确关系 \(AD\cdot BE=OF^{2}\).

半径垂直于切线，所以 \(OA\perp AD\)，\(OF\perp DE\)，\(OB\perp BE\). 从同一点向圆所作的两条切线长相等，故
\(DA=DF\)，\(EB=EF\).
考虑直角三角形 \(AOD\) 与 \(FOD\). 它们有公共斜边 \(OD\)，且 \(OA=OF\) 都是圆的半径，所以两三角形全等，从而 \(OD\) 平分 \(\angle AOF\)，即
\[
\angle AOD=\frac{1}{2}\angle AOF
\]
同理，直角三角形 \(BOE\) 与 \(FOE\) 全等，故 \(OE\) 平分 \(\angle BOF\)，即
\[
\angle BOE=\frac{1}{2}\angle BOF
\]
由于 \(A\)，\(O\)，\(B\) 共线，\(\angle AOF+\angle BOF=180^{\circ}\). 因此
\[
\angle BOE=90^{\circ}-\angle AOD
\]
又因为 \(\angle OBE=90^{\circ}\)，在直角三角形 \(BOE\) 中有
\[
\angle BEO=90^{\circ}-\angle BOE=\angle AOD
\]
所以直角三角形 \(AOD\) 与 \(BOE\) 相似，对应关系为 \(A\leftrightarrow B\)，\(O\leftrightarrow E\)，\(D\leftrightarrow O\). 于是
\[
\frac{AD}{BO}=\frac{AO}{BE}
\]
从而
\[
AD\cdot BE=AO\cdot BO=OF^{2}
\]
故由题面所给的 \(AB\cdot BE=OF^{2}\) 不能得到证明；能够由该构型恒成立并由原解答证明的是 \(AD\cdot BE=OF^{2}\).
\end{solution}

\begin{problem}
正六棱锥底面的边长是 \(4\mathrm{~cm}\)，侧面与底面所成二面角等于 \(45^{\circ}\)，求这棱锥体的体积.
\end{problem}

\begin{solution}
\solutionfigure{\bitmapfigure[max width=0.42\linewidth]{img/1958/hunan/q7_solution_fig1.png}}

设正六棱锥为 \(P\text{-}ABCDEF\)，底面中心为 \(O\)，高为 \(PO\). 在底面内过 \(O\) 作 \(OM\perp AB\)，垂足为 \(M\)，并连接 \(OA\)，\(OB\)，\(PM\). 因为 \(PO\perp\) 底面，且 \(OM\perp AB\)，由三垂线定理得 \(PM\perp AB\). 所以 \(\angle PMO\) 是侧面与底面所成二面角的平面角，从而
\[
\angle PMO=45^{\circ}
\]
正六边形可分成六个全等的正三角形，因此
\(OA=OB=AB=4\mathrm{~cm}\)，\(AM=MB=2\mathrm{~cm}\).
在直角三角形 \(BOM\) 中，由勾股定理
\[
OM=\sqrt{OB^{2}-MB^{2}}=\sqrt{4^{2}-2^{2}}=2\sqrt{3}\mathrm{~cm}
\]
又因为 \(PO\perp\) 底面，所以 \(PO\perp OM\)，三角形 \(POM\) 在 \(O\) 处为直角. 已知 \(\angle PMO=45^{\circ}\)，所以另一个锐角 \(\angle OPM=45^{\circ}\)，从而
\[
PO=OM=2\sqrt{3}\mathrm{~cm}
\]
底面正六边形的面积为六个三角形 \(AOB\) 的面积之和，即
\[
S_{\text{底}}=6\times\frac{1}{2}\times AB\times OM=6\times\frac{1}{2}\times4\times2\sqrt{3}=24\sqrt{3}\mathrm{~cm}^{2}
\]
所以棱锥体积为
\[
V=\frac{1}{3}S_{\text{底}}\cdot PO=\frac{1}{3}\times24\sqrt{3}\times2\sqrt{3}=48\mathrm{~cm}^{3}
\]
故棱锥体积为 \(48\mathrm{~cm}^{3}\).
\end{solution}

\begin{problem}
证明二面角的平分面内任意一点到这二面角的两个面的距离相等.
\end{problem}

\begin{solution}
\solutionfigure{\bitmapfigure[max width=0.42\linewidth]{img/1958/hunan/q8_solution_fig1.png}}

设二面角的两个面为平面 \(M\)，平面 \(N\)，棱为 \(BC\)，平分面为平面 \(P\). 在平面 \(P\) 内取任意一点 \(A\)，从 \(A\) 向平面 \(M\)，平面 \(N\) 分别作垂线，垂足为 \(D\)，\(E\)，于是
\(AD\perp\text{平面 }M\)，\(AE\perp\text{平面 }N\).
在平面 \(P\) 内过 \(A\) 作 \(AF\perp BC\)，其中 \(F\in BC\). 因为 \(BC\) 同时在平面 \(M\) 和平面 \(N\) 内，所以 \(AD\perp BC\)，\(AE\perp BC\). 过点 \(A\) 且含有直线 \(AD\) 和 \(AE\) 的截面平面垂直于 \(BC\)，它与 \(BC\) 的交点就是 \(F\)，故 \(A\)，\(D\)，\(E\)，\(F\) 共面.

在这个截面平面内，\(DF\) 是平面 \(M\) 的截线，\(EF\) 是平面 \(N\) 的截线，\(AF\) 是平面 \(P\) 的截线. 因为平面 \(P\) 是原二面角的平分面，所以 \(AF\) 平分截面内的角 \(DFE\)，从而
\[
\angle DFA=\angle AFE
\]
又因为 \(AD\perp\) 平面 \(M\)，\(DF\subset\) 平面 \(M\)，以及 \(AE\perp\) 平面 \(N\)，\(EF\subset\) 平面 \(N\)，所以
\[
\angle ADF=\angle AEF=90^{\circ}
\]
因此直角三角形 \(ADF\) 与 \(AEF\) 有公共斜边 \(AF\)，且有一个锐角相等，故两三角形全等. 对应的直角边相等，即
\[
AD=AE
\]
而 \(AD\)，\(AE\) 分别是点 \(A\) 到两个面的距离，所以二面角的平分面内任意一点到二面角两个面的距离相等.
\end{solution}
